\section{The clock of the frozen board}
\label{sec:clkfrozen}

The estimator is measured; point it at the real data.  This section
runs \eqref{eq:clkobj} on the two frozen ladders --- NVDA, where the
chapter began, and SPY --- and reads the results within the limits
\cref{sec:clkinverse} measured.

\Cref{fig:clkread}(a) is the run with a horizon.  Candidates are allowed
at the first five expiries --- through December 2026, a
\emph{horizon} choice recorded in the contract --- while the later
intervals carry no candidates and hold the ladder's far end in place
(they sit beyond the panel's right edge; panel~(b) shows them).  The orange steps are the calendar ladder of
\cref{fig:clkboard}(b); the blue steps are the same increments
divided by the solved clock.  The solver spends its budget exactly
where the chapter has been pointing: \MacClkReadHeroEarnD\ extra days
on the earnings-bearing interval (shaded), a combined
\MacClkReadHeroShortD\ days on the two short-dated intervals at the
left edge --- the board's hot Monday-to-Friday opening, a real if
modest acceleration --- and nothing anywhere else:
\MacClkReadHeroTotalD\ days in total.  Look at what the blue ladder
does.  (Units, once: both ladders are variance per year of their own
clock --- calendar year for orange, solved variance year for blue.)
The three intervals before the report meet at
\MacClkReadHeroFlatLevel\ --- the solver has equalized the
pre-earnings stretch --- and the earnings interval comes down from
\MacClkBoardHotFwd\ to \MacClkReadHeroEarnAfter, most of the
way to its neighbours but not all of it: the per-day charge keeps the
last few days unbought, the live shrinkage of
\cref{fig:clkident}(a) on a real board.  The December interval is
untouched --- orange peeking out behind blue --- and the in-horizon
spread of the ladder tightens from \MacClkReadHeroSpreadBeforeBp\ to
\MacClkReadHeroSpreadAfterBp\ variance bp.

Now read panel~(a) in two directions.  What the solver has
genuinely done is recover, \emph{from prices alone}, that the window
between August~21 and September~18 carries about seven ordinary days
more variance than its calendar length --- with no access to a news
feed.  The public schedule says why: NVDA's quarterly report sits in
that window.  But the solver does not know that, and cannot.  It
reads kinks; earnings, a product launch, a macro date, or a
genuinely lumpy term structure all look identical through
\eqref{eq:clkobj}.  The installed calendar is an \emph{estimate of
sizes at conventional dates}, to be reviewed against the public
schedule --- not an oracle that discovered earnings.

\begin{figure}[htbp]
\centering
\paperfig{\textwidth}{fig_clk_read.pdf}
\caption{Reading the frozen boards' clocks (forward variance per
calendar year, orange, vs per solved variance year, blue).
(a)~NVDA with candidates through year-end: \MacClkReadHeroEarnD\
extra days land on the earnings interval (shaded),
\MacClkReadHeroShortD\ on the short end, none elsewhere; the
pre-report intervals equalize at \MacClkReadHeroFlatLevel\ and the
in-horizon spread tightens from \MacClkReadHeroSpreadBeforeBp\ to
\MacClkReadHeroSpreadAfterBp\ variance bp.  (b)~NVDA with candidates
everywhere: \MacClkReadFullTotalD\ days installed, the ladder forced
flat (\MacClkReadFullSpreadBeforeBp\ to
\MacClkReadFullSpreadAfterBp\ variance bp) --- including
\MacClkReadFullMarchD\ days against a single ordinary interval.
(c)~SPY: \MacClkReadSpyTotalD\ day installed in total; the rising
ladder is left standing.}
\label{fig:clkread}
\end{figure}

Panel~(b) shows what happens when that review is skipped.  Same
board, same objective, but candidates at every expiry and no horizon.
The solver installs \MacClkReadFullTotalD\ extra days and flattens
the entire ladder, from a spread of \MacClkReadFullSpreadBeforeBp\
variance bp to \MacClkReadFullSpreadAfterBp\ --- including
\MacClkReadFullMarchD\ days against the December-to-March interval
alone, an ordinary three-month stretch whose elevated accrual is some
mixture of two more earnings reports and genuine term-structure
shape.  Nothing malfunctioned.  Handed enough candidates, the
flatness ideal is achievable, and the objective achieves it --- by
attributing \emph{everything} to the clock.  A perfectly flat
blue ladder is the optimizer satisfied, not the truth found; the
solver has no way to say ``this interval is hot for a reason the
clock should not absorb.''  The horizon of panel~(a), and the human
review of what it installs, are not decorations on the method ---
they are the method.

Panel~(c) is the control experiment the chapter owes the reader:
the index board.  SPY's calendar ladder rises steadily --- the
familiar upward term structure; its largest decrease anywhere is
\MacClkReadSpyDecBp\ variance bp --- and the solver, run with no
horizon at all, installs \MacClkReadSpyTotalD\ day in total, leaving
the ladder's \MacClkReadSpySpreadBp-variance-bp spread standing.
Both halves of the design show up in that restraint.  The
monotonicity charge in \eqref{eq:clkobj} fires only on
\emph{decreases} of forward variance, and this ladder has
essentially none; and at the index's forward volatilities --- the
board's median is \MacClkReadSpyMidVolPct\% --- the $\sigma^4$
materiality wall of \cref{sec:clkinverse} prices small events out
entirely.  A diversified index dilutes any
single name's report to a fraction of a day; the solver's silence on
SPY is the study's flat-ladder result repeated on real data.  The
single name keeps an event clock; the index keeps the calendar.

Three paragraphs of scope complete the chapter's subject.  First,
below one day.  This chapter's smallest unit of time is the day, and
that too is a convention: variance does not accrue evenly
\emph{within} a day either.  The classic evidence is the weekend ---
measured across decades, an exchange holiday adds far less variance
than a trading day, because most variance arrives while markets
trade \citep{FrenchRoll1986} --- and the same asymmetry separates a
session hour from an overnight gap.  The day-weight idea continues
below the day boundary without new structure: a day becomes a
profile (so much of its weight during the session, so much
overnight), and a scheduled morning print --- a CPI release at 8:30
--- is simply an event whose date $t_e$ is a fraction of a day,
entering \eqref{eq:clkclock} unchanged.  The reference implementation
carries such a sub-day clock; this book stops at day granularity,
stated here once.

Second, what the clock does not model.  The clock is a deterministic
forecast of the schedule of variance; it prices the \emph{anticipated}
part of an event and is silent on the realized surprise, on random
event sizes, and on volatility that is itself dynamic --- the
stochastic-volatility world, where an event is a random burst rather
than a known dilation.  That is model territory, developed in
\citet{Bergomi2016}, and it is the road this book declined at the end
of \cref{sec:clkboard}: richer, costlier, and unnecessary for the
readings this chapter set out to repair.

Third, where the idea comes from.  Reading asset prices on a
non-calendar clock is one of the oldest ideas in the subject:
\citet{Clark1973} explained heavy-tailed daily returns by
subordinating a Brownian motion to a transaction-volume clock, and
\citet{AneGeman2000} showed returns recover normality when read on
such an activity clock --- the empirical cousins of
\cref{thm:clkdds}.  This chapter's contribution to the reader is the
desk-sized version: a clock parameterized in days, priced into
readings by \cref{prop:clkgauge}, and estimable from the board with
stated limits.

\subsection{The contract}
\label{sec:clkcontract}

The chapter closes as the others do: three columns of decreasing
strength --- proved, measured, chosen.

\begin{table}[htbp]
\centering
\small
\caption{The chapter's contract.}
\label{tab:clkcontract}
\begin{tabular}{p{0.31\textwidth}p{0.31\textwidth}p{0.30\textwidth}}
\toprule
\textbf{Proved} & \textbf{Checked} & \textbf{Convention}\\
\midrule
The law at expiry knows the variance budget, not its schedule ---
exactly, for independent normal days; for every continuous martingale
by \cref{thm:clkdds} (cited). &
The walk of \cref{fig:clkwalk}: envelope kink
\MacClkWalkEnvKinkPct\% on the calendar axis, one smooth square root
on the variance axis. &
The clock $\tau$ is a deterministic forecast of the schedule;
realized surprises and random event sizes are out of scope
(\citealp{Bergomi2016}).\\
\addlinespace
Relabeling time (\cref{prop:clkgauge}): prices, butterfly validity
and calendar order invariant; readings covariant by
$\sqrt{t/\tau}$, interval rates by $\Delta t/\Delta\tau$. &
The worked week (\cref{fig:clkcrush}): ramp
\MacClkCrushRampStartPct\% to \MacClkCrushPeakPct\%, overnight crush
$-$\MacClkCrushDropBp\ vol bp, no unscheduled price move; the NVDA
dip (\MacClkBoardDipBp\ vol bp) diagnosed as the dodging expiry. &
Event currency: extra equivalent days $N_e$ per event, 365 days per
year; normalization \eqref{eq:clknorm} off throughout (a stated
budget choice).\\
\addlinespace
Linear-in-$\tau$ interpolation bends the dense curve at the
calendar's dates and nowhere else (\cref{sec:clkinterp}). &
The smear (\cref{fig:clkinterp}): $+\MacClkInterpOverBp$ vol bp
phantom before the event, $-\MacClkInterpUnderBp$ past it, zero at
every quote. &
Total variance linear in $\tau$ between quotes --- the
constant-accrual convention on the market's clock.\\
\addlinespace
Dates inside an interval are not identified
(\cref{prop:clkident}); an event can only lower its own interval's
$\mathcal{F}_i$. &
The study (\cref{fig:clkident}): shrinkage
\MacClkIdentShrinkStrongD/\MacClkIdentShrinkWeakD\,d at
$40\%/20\%$, quarterly wall at \MacClkIdentBlindD\,d, flat ladder
$\to$ \MacClkIdentFlatDays\ days. &
Candidate placement mid-interval; solver weights, half-day drop
threshold and horizon (appendix~\ref{app:clkprotocol}); solved sizes
are shrunk estimates.\\
\addlinespace
--- (this row is empirical by design) &
The frozen boards (\cref{fig:clkread}): \MacClkReadHeroEarnD\ days
onto NVDA's earnings interval (spread
\MacClkReadHeroSpreadBeforeBp$\to$\MacClkReadHeroSpreadAfterBp\ var
bp); the no-horizon run's \MacClkReadFullTotalD-day overreach; SPY
left at \MacClkReadSpyTotalD\ day. &
The installed calendar is reviewed against the public schedule
before use; the solver measures how much, never why.\\
\bottomrule
\end{tabular}
\end{table}

Step back to the book's thesis, one last time for Part~II.
\emph{Validity} was \cref{prop:clkgauge}: a clock is a relabeling,
and relabelings can never touch prices, densities, or order --- the
entire construction is arbitrage-proof by arithmetic.
\emph{Information} was what the board genuinely pins down: total
variances at quoted expiries, hence per-interval accrual rates ---
and, through the inverse problem, event \emph{sizes} above a
computable wall, but never dates inside an interval and never causes.
\emph{Convention} was everything else, each piece now on the table:
the day currency, the normalization choice, the interpolation rule,
the candidate placement, the solver's weights and horizon.  With this
chapter, Part~II's work is complete: the raw board has been turned
into forwards and discounts (Chapter~6), European prices
(Chapter~7), and now readings on the market's clock --- the observation
is fully prepared.  What Part~III asks is what none of this statics
can answer: when the spot moves, what happens to the smile?  The
surface gives every derivative in the strike direction and none in
the spot direction --- a missing derivative, and the next chapter is
about the conventions that supply it.
